Confidence in numerical assessments of coastal-defence performance under tsunami loading is strengthened by validation against independent observations of tsunami propagation and coastal impact.

Here, a one-way coupled 2D--3D framework is developed to reconstruct the 2011 Tohoku tsunami from coseismic generation to coastal impact.

Coseismic seafloor displacement defines the initial free-surface disturbance for a 2D nonlinear shallow-water model simulating propagation through a source-to-coast corridor.

Bathymetric and topographic datasets define the terrain, while absorbing layers and radiation conditions reduce spurious reflections at truncated open-ocean boundaries.

Extracted free-surface-elevation and depth-averaged velocity histories are mapped to time-dependent inlet conditions for a 3D unsteady Reynolds-averaged Navier--Stokes model with volume-of-fluid interface capture, simulating inundation and interaction with rigid defences and buildings.

Model performance is evaluated against observed arrival times, offshore waveforms, run-up heights and inundation extents using errors in timing, amplitude and spatial extent.

Expected outcomes are close quantitative agreement with event observations, physically consistent regional-to-local forcing, and resolved flow fields and hydrodynamic loads.

The resulting event-validated framework will enable comparative assessment of defence configurations under the historical event and systematically varied source scenarios.

Subsequent development will incorporate fluid--structure interaction and damage modelling to assess defence deformation, damage progression and failure margins under extreme tsunami loading.
